\section{Introduction}
The language of mathematical optimisation has become a fundamental tool for formulating an increasing number of real-world problems, where the aim is to maximise or minimise an objective function, often subject to constraints. Whenever the problem's objective and/or the constraints depend on a set of parameters, $x$, representing variables that do not change the problem's structure but produce a specific problem instantiation, it takes the form of a parametric optimisation program. 

Formally, we represent the decision variable as $z \in \mathbb{R}^n$, the parameter variable as $x \in \mathcal{X} \subset \mathbb{R}^d$, drawn from parameter space $\mathcal{X}$, the objective function as $f : \mathbb{R}^n\times\mathbb{R}^d \rightarrow \mathbb{R}$, and a constraint set as $C \subseteq \mathbb{R}^n$. The general constrained parametric optimisation problem $P(x)$ can then be formulated as:

\begin{equation} \label{eq:parametric_optimisation}
\begin{array}{ll}
    \text{minimise}  & f(z,x) \\
    \text{subject to} & z \in C(x).
\end{array}
\end{equation}

The formulation in~\eqref{eq:parametric_optimisation} can be deployed in a vast number of practical applications, including (but not limited to) portfolio optimisation in finance, hyperparameter tuning in machine learning (ML) applications such as support vector machines (SVMs) or Lasso regression, and model predictive control (MPC) in control engineering. In real-time and safety-critical domains, such as driver-assistance systems, autonomous robotics, medical imaging pipelines, power grid dispatch, and aerospace guidance, it is crucial to instantiate and solve $P(x)$ rapidly and predictably, as failures can compromise finite-iteration safety or stability. In general, solutions to problem $P(x)$, denoted as $z^\star(x) \defeq \arg\min_{z \in C(x)} f(z,x)$, carry a dependency on $x$ and therefore constitute a \textit{solution map}, which can be exploited to evaluate optimal decisions through a simple function evaluation, completely bypassing the need to re-solve~\eqref{eq:parametric_optimisation} when a new (potentially unseen) $x$ is presented. This structural feature is unique to parametric programs, and is critical for applications that demand real-time or safety-critical decision-making.

We refer to any mechanism to compute $z^\star(x)$ from an initial condition $z^0 \in Z$ as an \emph{algorithm}. A popular family of algorithms used to solve several classes of optimisation problems is that of \textit{first-order methods}: iterative schemes relying only on first-order derivative information (e.g., gradients or subgradients) of the functions involved. 
These are typically denoted as:
\begin{equation}\label{eq:first_order_method}
  z^{k+1}(x) = T(z^k(x), x),\quad k=0,1,\dots,K
\end{equation}
where $T: \mathbb{R}^n \times \mathbb{R}^d \rightarrow \mathbb{R}^n$ gives the gradient-dependent mapping\footnote{
As an example, consider the \textit{projected gradient descent} algorithm, given by $T(z^k(x), x) \defeq \Pi_{C(x)}(z^k-\theta\nabla_zf(z^k,x))$, where $\nabla_zf(z^k,x)$ is the gradient in the direction of $z$, and $\theta$ is a hyperparameter representing the step size} 
from $z^k$ to $z^{k+1}$. First-order methods offer two main advantages over other algorithm types: they are extremely cheap to implement and scale, and they are robust to the choice of initial conditions set $Z$ (also called a \textit{warm-start} set). Their usage has seen a recent resurgence in applications for both microchip-scale embedded optimisation and high-compute large-scale optimisation~\cite{StellatoPresentation2025}. They have been applied to linear (Google PDLP~\cite{googlePDLP}, NVIDIA cuOPT~\cite{NVIDIAcuopt}), quadratic (OSQP~\cite{stellato2020osqp}, ProxSuite~\cite{proxsuite}), and conic (SCS~\cite{SCSconic}) programs. They have also become ubiquitous in the ML literature~\cite{ruder2016overview}. However, they can be found to converge too slowly for certain problem instances, which represents a major setback in safety-critical applications where only a specific number of iterations of algorithm~\eqref{eq:first_order_method} can be afforded. A fundamental challenge is therefore to provide performance guarantees on the convergence behaviour of first-order methods to ensure their feasibility in practical implementations. This will be the principal focus of this Senior Thesis project.

Performance analysis strategies to provide convergence guarantees for first-order methods can be broadly split into two categories: classical and computer-assisted. Classical techniques derive analytical convergence bounds through mathematical proofs, often relying on problem-specific assumptions like smoothness or convexity to establish worst-case rates. Computer-assisted techniques, by contrast, leverage computational tools like semidefinite programming and performance estimation problems (PEPs) to numerically compute tight worst-case bounds, enabling automated verification and optimisation of algorithm performance without exhaustive manual analysis. In general, these study some performance metric $\phi(z,x)$, such as suboptimality $f(z^k(x),x) - f(z^\star(x),x)$, distance to optimality $\|z^k(x) - z^\star(x)\|$, or distance travelled $\|z^{k+1}(x) - z^k(x)\|$, for some choice of norm $\|\cdot\|$.
Guarantees are then framed either in a probabilistic sense, i.e., bounding:

\begin{equation}\label{eq:probabilistic_guarantee}
\mathbb{P}\{\phi(z^K(x),x) > \varepsilon\}
\end{equation}

for some probability measure $\mathbb{P}$ and tolerance $\varepsilon$, or in a worst-case sense, i.e., evaluating and bounding the worst-case performance $\phi_{\max}(z^K(x),x) \defeq \max_{x\in\mathcal{X}} \phi(z^K(x),x)$. Recent work has given particular attention to worst-case analysis: for instance, in \cite{VRBS25}, the authors investigated finite-step convergence verification of first-order methods for parametric convex quadratic programs, obtaining tight, data-free quantifications of $\phi_{\max}(z^K(x),x)$ that incorporate warm-starting effects; similarly, \cite{VRBS+24} formulated an exact mixed-integer linear program to maximise $\phi(z^K(x),x)$, defined using the $\infty$-norm, by encoding the algorithm’s dynamics as polyhedral constraints. While these approaches yield rigorous guarantees, they tend to be conservative for typical parameter instances. To address these limitations, probabilistic analyses have emerged as a promising alternative: for instance, \cite{RSBS25} developed probabilistic guarantees for both classical and learned optimisers using sample convergence bounds and the Probably Approximately Correct (PAC)-Bayes framework, respectively, while \cite{JHKM+25} employed advancements in the so-called Scenario Approach Theory to generate tight, data-driven guarantees for both convergence and general definitions of $\phi(z^K(x),x)$. However, these approaches suffer from important drawbacks: the PAC-Bayes framework is unintuitive and potentially limited in proving generalisation for certain learnable tasks, whereas Scenario Approach bounds are heavily sample-count dependent, which can lead to overly conservative guarantees with insufficient data.